\chapter{Ingrown Hair}
\index{Ingrown Hair}

\section*{Definition and Overview}
An ingrown hair is a hair that curls back and grows into the skin instead of emerging through the surface. It typically causes a small red, tender bump, often with a visible hair beneath the skin, and commonly affects areas that are shaved or waxed, such as the face, neck, legs, and bikini area. Ingrown hairs are usually harmless and resolve on their own, but they can become infected, form pus-filled lesions, and cause scarring and skin discolouration, particularly in people with curly hair. Prevention centres on shaving technique, and treatment ranges from simple home care to removal of the embedded hair.

\section*{Epidemiology}
Ingrown hairs are very common in people who shave regularly, and the risk is highest in people with naturally curly hair, because the curved hair shaft is more likely to re-enter the skin. Men who shave the face and neck, and people who shave or wax the legs, bikini area, or underarms, are most affected. The condition is more common with frequent shaving, shaving against the grain, tight clothing, and the use of dull blades. While almost everyone will have an ingrown hair at some point, some people develop them recurrently with irritation and scarring.

\section*{Aetiology and Pathophysiology}
When a hair is cut short, the sharpened tip can fail to penetrate the skin surface as it grows. Instead, it curls back and grows into the epidermis or dermis, where the body mounts an inflammatory response, producing a papule or pustule. Shaving against the direction of hair growth leaves shorter, sharper hairs; tight skin during shaving allows cut hairs to retract below the surface; and curly hair is geometrically prone to re-entry. Blocked hair follicles can become infected with skin bacteria, most commonly Staphylococcus, turning a simple ingrown hair into folliculitis or a boil. Chronic cases may cause post-inflammatory hyperpigmentation and scarring.

\section*{Clinical Features}
\begin{itemize}
    \item Small red bumps or pustules at the site of shaving or waxing
    \item A visible hair under the skin surface
    \item Itching, tenderness, or pain in the affected area
    \item Pus or crusting when follicles become infected
    \item Darkening of the skin and scarring with repeated episodes
    \item Common sites: beard area, neck, legs, bikini line, and armpits
    \item Recurrence at the same sites with continued shaving
\end{itemize}

\section*{Differential Diagnosis}
Other skin conditions can resemble ingrown hairs.
\begin{itemize}
    \item \textbf{Folliculitis:} infection of the hair follicle, sometimes without an ingrown hair
    \item \textbf{Bacterial boils and abscesses:} deeper, larger, more painful infections
    \item \textbf{Acne:} comedones and inflamed lesions, often on the face
    \item \textbf{Contact dermatitis:} irritation or allergy to shaving products
    \item \textbf{Skin infections:} impetigo and fungal infections
    \item \textbf{Molluscum contagiosum and warts:} skin lesions that can be mistaken for ingrown hairs
    \item \textbf{Rare skin tumours:} persistent or unusual lesions require assessment
\end{itemize}

\section*{When to Seek Medical Care}
Most ingrown hairs settle with home care. Medical review is appropriate for lesions that are painful, enlarging, recurrent, or showing signs of infection such as spreading redness, pus, fever, or increasing tenderness. People with diabetes or weakened immunity should seek earlier assessment. A persistent or unusual lesion that does not heal should be checked.

\textbf{Contact your local emergency services for:}
\begin{itemize}
    \item Rapidly spreading redness, pain, and swelling, which may indicate a serious skin infection
    \item Fever, chills, or feeling generally unwell with a skin lesion
    \item Red streaks extending from the lesion, which may indicate lymphangitis
\end{itemize}

\section*{Diagnosis and Investigations}
\begin{itemize}
    \item \textbf{History and examination:} the site, shaving routine, and appearance of the lesions
    \item \textbf{Visual inspection:} identification of the embedded hair and surrounding inflammation
    \item \textbf{Microbiological swab:} if infection is suspected, to identify bacteria
    \item \textbf{Skin biopsy:} rarely needed, for persistent or atypical lesions
    \item \textbf{Assessment of risk factors:} hair type, shaving methods, and skin care products
\end{itemize}

\section*{Management}
\begin{itemize}
    \item \textbf{Warm compresses:} applied for a few minutes several times a day to soften the skin and help the hair emerge
    \item \textbf{Gentle release:} a sterile needle or tweezers to lift the hair loop above the skin surface, without digging for deeply buried hairs
    \item \textbf{Stop shaving:} temporarily allowing the hairs to grow out
    \item \textbf{Improved shaving technique:} shaving in the direction of hair growth with a sharp blade and warm water
    \item \textbf{Exfoliation:} gentle physical or chemical exfoliation with products such as salicylic acid
    \item \textbf{Antibacterial washes:} for recurrent folliculitis
    \item \textbf{Alternatives to shaving:} laser hair removal or electrolysis for persistent problems
    \item \textbf{Treatment of infection:} antibiotics by mouth or cream where bacterial infection is present
\end{itemize}

\section*{Complications}
\begin{itemize}
    \item Infection of the follicle, progressing to boils and abscesses
    \item Cellulitis in severe or neglected cases
    \item Scarring and permanent skin darkening (hyperpigmentation)
    \item Chronic irritation and pain, especially in the beard area
    \item Rarely, the formation of a sinus tract or keloid
    \item Psychological effects of visible, recurrent lesions
\end{itemize}

\section*{Prognosis and Outcomes}
Most ingrown hairs resolve within days to weeks with simple measures, and releasing the trapped hair gives rapid relief. With consistent changes to shaving technique and skin care, recurrence is markedly reduced. For people with severe or recurrent problems, laser hair removal offers the best long-term results. Complications such as infection and scarring are usually preventable with early treatment and appropriate technique.

\section*{Prevention and Self-Care}
\begin{itemize}
    \item Shave in the direction of hair growth with a sharp, clean blade
    \item Prepare the skin with warm water and shaving cream
    \item Do not stretch the skin tightly while shaving
    \item Use a single-blade razor or an electric razor if prone to ingrown hairs
    \item Exfoliate gently a few times a week
    \item Apply moisturiser or aftershave designed for sensitive skin
    \item Avoid shaving too closely or too frequently
    \item Do not pick at or dig out ingrown hairs
    \item Consider laser hair removal for recurrent problems
\end{itemize}

\section*{Sources and Further Reading}
The following reputable sources provide further information for healthcare students and patients seeking reliable, current advice about ingrown hairs.
\begin{itemize}
    \item Healthdirect Australia. \textit{Ingrown hairs.} https://www.healthdirect.gov.au/
    \item Better Health Channel. \textit{Ingrown hairs and shaving.} https://www.betterhealth.vic.gov.au/
    \item Australasian College of Dermatologists. \textit{Patient information on skin conditions.} https://www.dermcoll.edu.au/
    \item NHS. \textit{Ingrown hairs.} https://www.nhs.uk/conditions/ingrown-hairs/
    \item Mayo Clinic. \textit{Ingrown hairs.} https://www.mayoclinic.org/
    \item American Academy of Dermatology. \textit{How to prevent ingrown hairs.} https://www.aad.org/
\end{itemize}
